\documentclass{article}
\title{Notas de Aula - MAT01059}
\author{Leonardo Suarez}
\date{\today}
\usepackage{macros}

\begin{document}
\maketitle

Notas de aula de Análise Matemática C.

\tableofcontents

\section{Funções Reais de \(n\) Variáveis}
\subsection{Derivadas parciais}
\subsection{Funções de classe \(C^1\)}
\subsection{O Teorema de Schwarz}
\subsection{A fórmula de Taylor}
\subsection{Pontos críticos}
\subsection{Funções convexas}

\section{Funções Implícitas}
\subsection{Uma função implícita}
\subsection{Hiperfícies}
\subsection{Multiplicador de Lagrange}

\section{Aplicações Inversas e Implícitas}
\subsection{O Teorema da Aplicação Inversa}
\subsection{Várias Funções Implícitas}

\section{Superfícies Diferenciáveis}
\subsection{Parametrizações}
\subsection{Superfícies diferenciáveis}
\subsection{O espaço vetorial tangente}
\subsection{Superfícies orientáveis}
\subsection{Multiplicadores de Lagrange}
\subsection{Aplicações diferenciáveis entre superfícies}

\section{Integrais Múltiplas}
\subsection{A definição de integral}
\subsection{Conjuntos de medida nula}
\subsection{Cálculo com integrais}
\subsection{Conjuntos \(J\)-mensuráveis}
\subsection{A integral como limite de somas de Riemann}

\section{Mudança de Variáveis}
\subsection{O caso unidimensional}
\subsection{Difeomorfismos primitivos}
\subsection{Todo difeomorfismo \(C^1\) é localmente admissível}

\section{Formas Alternadas}

\section{Formas Diferenciais}

\section{O Teorema de Stokes}

\end{document}